\documentclass{article}

\usepackage{corl_2026}

\title{基于目标流优化的固定翼飞行技能组合方法}

\author{
  作者一\\
  单位\\
  \texttt{email} \\
  \And
  作者二\\
  单位\\
  \texttt{email} \\
}

\begin{document}
\maketitle

%===============================================================================

\begin{abstract}
  直接从作动器指令学习敏捷的固定翼机动飞行仍然面临挑战，原因包括姿态--速度--能量动力学的强耦合、欧拉角奇异性问题，以及仅用位置误差评估环形机动的局限性。
  本文提出一个组合式框架，其核心是一个冻结的、无 PID 的、基于四元数的强化学习（RL）飞行技能，该技能在跟踪航向、俯仰、滚转和速度（H/P/R/Vt）目标流的同时直接输出作动器指令。
  不同于使用固定启发式方法生成该目标流，我们引入了\emph{滚动时域目标流优化}（Receding-Horizon Target-Stream Optimization, RH-TSO）：在每个决策步，候选目标流参数化方案通过冻结飞行技能的闭环 rollout 进行评估，并根据几何、能量、安全性和平滑性代价进行选择。
  优化在可执行的目标流空间（前瞻距离、目标速度、俯仰塑形）中进行，而非直接在作动器空间中搜索，使得搜索紧凑且接口可解释。
  我们进一步应用能量感知微调课程，将四元数技能扩展到垂直弧线机动。
  在高保真 F-16 仿真中，我们证明：（1）对于超过 60° 俯仰的机动，四元数目标编码显著优于欧拉角编码；（2）RH-TSO 在曲线和垂直机动上相比固定目标流参数将跟踪精度提高 24--47\%，且不存在适用于所有机动类别的单一固定参数化方案；（3）能量感知微调使垂直弧线最高可达 150°，达到 B 级环面质量，而 180° 半筋斗仍为已诊断的失败边界；（4）传统的横偏误差（CTE）单独评估对环形机动具有误导性——速度切线误差、机头切线误差、翼面误差和四元数测地误差等几何感知指标能更忠实地反映机动质量。
\end{abstract}

\keywords{固定翼机动飞行，强化学习，四元数控制，目标流优化，几何感知评估}

%===============================================================================

\section{引言}

直接从作动器指令学习敏捷的固定翼飞行，本质上比研究成熟的四旋翼问题更困难。
固定翼飞机将姿态、速度、能量和翼面几何耦合在一起：滚转指令改变升力方向从而影响转弯平面；俯仰变化将动能转化为高度；每个机动都必须满足失速限制、结构过载约束以及气动舵面的物理带宽。
这些约束使端到端学习复杂机动变得困难，也使\emph{评估}机动是否被正确执行同样困难。

常见的做法是将问题分解为：经典内环控制器（PID 或自动驾驶仪）跟踪姿态和速度指令，RL 策略在更高层面输出这些指令。
虽然实际有效，但这种分解将系统与手工设计的内环增益绑定，这些增益可能无法在整个飞行包线内通用；且通常使用欧拉角表示，在俯仰角接近 $\pm90^\circ$ 时出现万向节死锁——而这恰恰是垂直机动运行的空域。

本文采用一种不同的方法。
我们训练了一个\emph{无 PID}、基于四元数的 RL 飞行技能，该技能直接输出作动器指令（油门、升降舵、副翼、方向舵、减速板），同时以一组目标航向、俯仰、滚转和速度值为条件。
该技能在随机化的 H/P/R/Vt 目标上一次训练完成后即被冻结。
关键问题随之变为：\emph{我们应该向这个冻结的技能输入怎样的目标流，才能执行期望的机动？}

一个简单的回答是启发式地固定目标流参数——例如，恒定的前瞻距离和恒定的目标速度。
我们的实验表明，这样做会失败：水平曲线任务（S 形转弯、8 字形）偏好较短的前瞻距离和较低的目标速度，而垂直拉起任务则偏好较长的前瞻距离和较高的目标速度。
不存在适用于所有机动类别的单一固定参数化方案。

因此，我们提出\textbf{滚动时域目标流优化（RH-TSO）}：在每个决策步，我们生成候选目标流参数化方案，通过冻结飞行技能的闭环 rollout 评估每个方案，并选择使几何、能量、安全性和平滑性组合代价最小化的方案。
优化在紧凑的目标流参数空间（前瞻距离、目标速度、俯仰塑形偏移）中进行，而非在维数极高的作动器空间中。
这种设计使搜索保持可追踪，同时尊重技能的执行约束：每个候选方案都由即将执行它的同一冻结策略评估，因此不可行的目标会自然地受到惩罚。

本文做出以下贡献：
\begin{enumerate}
  \item 一个\textbf{无 PID 的四元数 RL 飞行技能}，直接输出作动器指令，以动态 H/P/R/Vt 目标为条件，避免了欧拉角奇异性。
  \item \textbf{滚动时域目标流优化（RH-TSO）}，通过确定性并行打靶搜索配合闭环策略 rollout，在可执行技能空间中优化目标流参数，替代固定启发式参数化方案。
  \item 一个\textbf{能量感知微调课程}，通过几何塑形奖励将四元数技能扩展到垂直弧线机动（60°--150°），同时保留水平飞行性能。
  \item 一个\textbf{几何感知评估协议}，证明 CTE 单独评估环形机动可能产生虚假的成功判断，并提出一套从姿态、速度、能量和几何维度综合评估机动质量的多指标体系。
\end{enumerate}

我们在高保真 F-16 仿真中验证了所有组件，使用 NASA TP-1538 气动系数表和 6 自由度非线性动力学，仿真频率 50 Hz。

%===============================================================================

\section{相关工作}

我们将相关工作分为四组。

\subsection{空中机器人学习与敏捷飞行}

大多数 CoRL 敏捷飞行论文以四旋翼为研究对象。
Deep Drone Racing~\citep{kaufmann2023deep} 展示了学习到的感知到航点流水线与经典规划器和控制器相结合可以实现竞速级别的性能。
DATT~\citep{huang2023datt} 展示了学习到的轨迹跟踪策略可以在气动干扰下处理激进的飞行轨迹。
Flightmare~\citep{song2021flightmare} 为 RL 原型开发提供了灵活的四旋翼仿真基础设施。

我们的工作研究了一个互补问题：固定翼机动组合，其中姿态、速度、能量和翼面几何以四旋翼动力学无法捕捉的方式紧密耦合。
固定翼飞机无法悬停，不能独立控制姿态和速度，必须在整个机动过程中管理失速裕度和结构过载。

\subsection{学习与规划/控制的接口}

机器人学习中一个反复出现的主题是：中间可执行表示通常优于从任务规范直接映射到作动器的端到端方法。
Deep Drone Racing 使用航点和期望速度作为中间表示~\citep{kaufmann2023deep}。
将最优控制与学习相结合的工作~\citep{kaufmann2019combining} 表明，学习到的感知模块配合最优控制器在新环境中优于纯端到端学习。
实时生成时间最优四旋翼轨迹的工作~\citep{penicka2022real} 证明了在学习到的潜空间中进行轨迹优化的可行性。

我们的中间表示——H/P/R/Vt 目标流——是为固定翼领域专门设计的：它编码完整的姿态和速度信息，避免欧拉奇异性，并使冻结的底层技能能够泛化到新机动。
与使用固定规划器的先前工作不同，我们通过冻结技能自身的闭环 rollout \emph{在线优化}目标流参数。

\subsection{直接学习控制与 RL 飞行技能}

四旋翼集群的分散控制~\citep{batra2022decentralized} 证明了端到端 RL 可以在无手工设计内环的情况下实现编队飞行。
软体多旋翼控制~\citep{song2021soft} 使用神经网络动力学辨识为经典 PID 调节不切实际的系统学习控制策略。
DATT 展示了在动态不确定性下直接策略学习用于轨迹跟踪~\citep{huang2023datt}。

与依赖 PID/自动驾驶仪内环的高层 RL 方法不同，我们的底层技能直接输出固定翼作动器指令，同时以四元数编码的目标姿态和速度为条件。
该策略是无 PID 的，学习了从状态到作动器的完整映射，包括隐式的姿态角速率阻尼和转弯协调。

\subsection{超越位置跟踪的评估}

大多数轨迹跟踪论文仅用位置误差评估。
DATT 认识到这一不足，引入了姿态跟踪和鲁棒性指标~\citep{huang2023datt}。
对于固定翼环形机动，位置误差本身的误导性更大：一架飞机可以紧贴参考轨迹飞行，但偏航方向任意、侧滑角很大、或翼面几何倒置——这些都表现为"低 CTE"但代表了根本错误的机动执行。

我们使用速度切线误差、机头切线误差、翼面误差和四元数测地误差，以及领域特定的安全指标（攻角裕度、过载、能量保持、空速包线）进行评估。

%===============================================================================

\section{方法}

我们的系统由三个组件构成：
(1)~一个基于四元数的 RL 飞行技能，在给定状态和 H/P/R/Vt 目标观测的条件下直接输出作动器指令；
(2)~滚动时域目标流优化（RH-TSO），通过闭环策略 rollout 评估候选方案来选择目标流参数；
(3)~一个能量感知训练课程，将飞行技能扩展到垂直机动。

% 图1：系统概览
% TODO: 参考机动 → RH-TSO (目标流优化循环) → [H,P,R,Vt]流 → 冻结的四元数RL技能 → 作动器指令 → 高保真仿真 → 几何感知代价 → RH-TSO反馈

\subsection{基于四元数的 RL 飞行技能}

底层飞行技能通过 PPO 算法训练，使用基于 GRU 的 Actor-Critic 架构。
它接收 21 维观测向量，在五个控制通道上输出离散动作：油门（31 档）、升降舵（41 档）、副翼（41 档）、方向舵（41 档）和减速板（5 档）。
该策略完全无 PID：不使用任何经典内环控制器。

\subsubsection{观测空间}

观测向量是无奇异的：

\begin{itemize}
  \item \textbf{四元数误差向量部}（3 维）：$q_{\text{err}} = q_{\text{tgt}} \otimes q_{\text{curr}}^*$ 的向量分量，采用最短旋转约定（$w \geq 0$）。
  \item \textbf{归一化速度误差}（1 维）：$(v_t - v_{t,\text{tgt}}) / 340$。
  \item \textbf{归一化高度}（1 维）：$h / 5000$ m。
  \item \textbf{归一化速度}（1 维）：$v_t / 340$。
  \item \textbf{机体坐标系中的目标方向}（3 维）：$[\cos\theta\cos\psi, \cos\theta\sin\psi, \sin\theta]$ 旋转到体轴系中。
  \item \textbf{角速率}（3 维）：$P, Q, R$。
  \item \textbf{攻角和侧滑角}（4 维）：$\sin\alpha, \cos\alpha, \sin\beta, \cos\beta$。
  \item \textbf{历史动作}（5 维）：归一化的上一步指令，用于时间平滑性。
\end{itemize}

所有值被裁剪到有界范围内。四元数公式消除了在 $\pm90^\circ$ 俯仰角处的欧拉奇异性。

\subsubsection{奖励函数}

基准奖励是姿态和速度跟踪的加权几何平均：

\begin{equation}
  r_{\text{base}} = (r_{\text{att}})^{0.8} \cdot (r_{\text{speed}})^{0.2}
\end{equation}

\begin{equation}
  r_{\text{att}} = \exp\left(-\left(\frac{\theta}{\theta_{\text{scale}}}\right)^2\right), \quad \theta = 2\arccos(|w_{\text{err}}|), \;\; \theta_{\text{scale}} = 5^\circ
\end{equation}

\begin{equation}
  r_{\text{speed}} = \exp\left(-\left(\frac{v_t - v_{t,\text{tgt}}}{24~\text{m/s}}\right)^2\right)
\end{equation}

基准奖励辅以高度安全项。总奖励裁剪到 $[0, 1]$，由智能体存活状态掩码。

\subsubsection{离散动作解码}

离散动作通过等距分档解码为连续指令：

\begin{align}
  \text{油门} &\in [0, 1] \quad \text{(31 档)} \\
  \text{升降舵} &\in [-1, 1] \quad \text{(41 档)} \\
  \text{副翼}  &\in [-1, 1] \quad \text{(41 档)} \\
  \text{方向舵}   &\in [-1, 1] \quad \text{(41 档)} \\
  \text{减速板} &\in [0, 1] \quad \text{(5 档)}
\end{align}

减速板头以 $[0, -1.5, -1.5, -1.5, -1.5]$ 的偏置初始化，默认为收起状态。

\subsection{目标流公式化}

RL 飞行技能的输入不是航点（空间位置），而是一个连续的\emph{动态姿态和速度目标}流 $\tau_t = (\psi_t, \theta_t, \phi_t, v_t)$。
对于给定的参考机动，目标流生成器 $G$ 以参数向量 $\bm{\lambda}$ 产生该目标流：

\begin{equation}
  \{\tau_t\}_{t=0}^{T} = G(\text{机动}, \bm{\lambda})
\end{equation}

参数向量 $\bm{\lambda}$ 控制参考机动如何转换为 H/P/R/Vt 目标。具体而言，在我们的实现中：

\begin{itemize}
  \item $\lambda_{\text{lookahead}}$：沿轨迹的前瞻距离（m），决定哪个未来参考点提供目标航向和俯仰。
  \item $\lambda_{v_t}$：整个机动过程中指令的目标速度（m/s）。
  \item $\lambda_{\text{pitch\_offset}}$：作用于参考俯仰的加性偏移（rad），用于垂直塑形。
\end{itemize}

关键洞见是：$\bm{\lambda}$ 位于一个低维空间（通常 2--3 个参数），但不同选择会因机动类别不同而产生显著不同的闭环跟踪质量。

\subsection{滚动时域目标流优化（RH-TSO）}

给定冻结的飞行技能 $\pi$、一个参考机动和当前飞机状态，RH-TSO 在每个决策步通过求解以下问题来选择目标流参数：

\begin{equation}
  \bm{\lambda}^* = \arg\min_{\bm{\lambda} \in \Lambda} \; J\big(\text{rollout}(\pi, G(\text{机动}, \bm{\lambda}), s_0), \text{ref}\big)
\end{equation}

其中 $J$ 是组合了几何误差、能量偏差、安全违规和控制平滑性的代价函数，$\text{rollout}(\cdot)$ 在有限时域 $H$ 上执行冻结策略跟踪候选目标流的闭环仿真。

\subsubsection{搜索策略}

对于当前低维参数空间，我们使用\textbf{确定性并行打靶格点搜索}（deterministic parallel-shooting lattice search）：候选 $\bm{\lambda}$ 值的网格通过冻结策略的批量闭环 rollout 并行评估。
每个候选方案产生一条轨迹；使 $J$ 最小化的方案被选中，其第一步动作被应用。
然后优化窗口前移一步，重复上述过程。

这不是对任意参数的枚举——每个候选目标流都\emph{通过即将执行它的同一冻结策略}进行评估，因此搜索被约束在学习技能的可执行流形上。
不可行或动态不一致的目标自然会获得高代价，因为策略无法很好地跟踪它们。

\subsubsection{代价函数}

RH-TSO 代价是时域 $H$ 上的加权和：

\begin{equation}
  J = w_g J_{\text{geo}} + w_e J_{\text{energy}} + w_s J_{\text{safety}} + w_u J_{\text{smooth}}
\end{equation}

其中：
\begin{itemize}
  \item $J_{\text{geo}}$：平均横偏误差和相对参考的姿态偏差。
  \item $J_{\text{energy}}$：相对参考比能 $E/m = \frac{1}{2}v_t^2 + gh$ 的偏差。
  \item $J_{\text{safety}}$：失速（$v_t < 150$ m/s）、超 G（$>9$ G）和触地惩罚。
  \item $J_{\text{smooth}}$：作动器变化率惩罚 $\sum \|u_t - u_{t-1}\|^2$。
\end{itemize}

\subsubsection{向高维空间的扩展}

当目标流参数空间增大时（如时变速度剖面、分段前瞻距离等），格点搜索变得昂贵。
在此情况下，基于采样的方法如交叉熵方法（CEM）或模型预测路径积分（MPPI）控制可以替代格点搜索，同时保持相同的结构：候选方案通过闭环策略 rollout 评估，代价梯度从轨迹代价中估计。
我们将其作为自然扩展加以说明，但不在当前实验中声称任何特定搜索方法的优越性。

\subsection{能量感知垂直扩展}

基准四元数技能在随机化 H/P/R/Vt 目标上训练，对最高约 $\pm60^\circ$ 的俯仰角能实现可靠跟踪。
超过此范围，两种失效模式出现：（1）持续拉起过程中的能量消耗导致失速；（2）环面漂移导致不协调飞行。

我们应用一个能量感知微调课程来扩展奖励函数，同时混入原始任务回放以防止灾难性遗忘。

\subsubsection{课程设计}

由 \texttt{vertical\_stage} 计数器控制的分阶段难度递进：

\begin{itemize}
  \item \textbf{阶段 0--1}：渐变俯仰目标（$5^\circ$--$10^\circ$），中等半径（8000--10000 m），定高代理任务。
  \item \textbf{阶段 2--3}：较大俯仰目标（$15^\circ$--$20^\circ$），60° 垂直弧线，紧凑半径圆形代理任务。
  \item \textbf{阶段 4--7}：拉起机动，90° 垂直弧线，半筋斗保持阶段（半径 2000--5000 m）。
\end{itemize}

原始任务混合概率在 300 次任务切换中从 0.80 退火到 0.05，以保留水平飞行能力。

\subsubsection{能量感知奖励塑形}

微调奖励在基准奖励基础上增加：

\begin{equation}
  r = 0.68 r_{\text{att}} + 0.24 r_{\text{speed}} + r_{\text{低速}} + r_{\text{能量}} + r_{\text{爬升}} + r_{\text{定高}} + r_{\alpha\beta} + r_{G} + r_{\text{平滑}} + r_{\text{存活}}
\end{equation}

关键新增项：低于 180/170 m/s 的非对称低速惩罚；每步和每回合的能量保持惩罚；由速度安全门控的爬升进度奖励；攻角/侧滑软惩罚（$>15^\circ$/$>10^\circ$）；过载惩罚（$>9$ G，硬上限 15 G）；动作平滑性；以及针对半筋斗机动的环面几何塑形（机头切线、翼面、速度切线、机头速度对齐），由倒飞姿态检测器门控。

\subsection{训练过程}

基准技能使用 PPO 训练 600 个 epoch（512 并行环境 $\times$ 512 步，GRU 隐藏维 256，Adam 学习率 $2.5\times10^{-4}$）。
能量感知微调从此检查点开始，使用扩展课程额外训练 619 个 epoch。
RH-TSO rollout 中使用的冻结技能是微调后的检查点。

%===============================================================================

\section{实验设置}

\subsection{仿真平台}

所有实验使用高保真 F-16 仿真，实现 NASA TP-1538 气动系数表，采用基于 JAX 的三线性插值，6 自由度非线性动力学，仿真频率 50 Hz，RL 智能体交互频率 5 Hz。
前缘襟翼（LEF）按 JSBSim 约定基于攻角和马赫数自动调度。

\subsection{基线方法与变体}

\begin{itemize}
  \item \textbf{欧拉基线}：使用欧拉角观测和奖励的 PPO 策略。
  \item \textbf{四元数（本文方法）}：基于四元数的观测和奖励（第 3.1 节）。
  \item \textbf{四元数 + 垂直能量（本文方法）}：从四元数检查点进行能量感知微调。
  \item \textbf{RH-TSO（本文方法）}：冻结的四元数+VE 策略配合滚动时域目标流优化。
  \item \textbf{固定 $\bm{\lambda}$}：冻结的四元数+VE 策略配合单一固定目标流参数化（RH-TSO 对比基线）。
\end{itemize}

\subsection{机动测试集}

\begin{itemize}
  \item \textbf{S 形转弯}：正弦航向 $\pm45^\circ$，周期 85 s。
  \item \textbf{8 字形}：倍频调制航向，周期 120 s。
  \item \textbf{螺旋线}：组合正弦航向和俯仰变化。
  \item \textbf{水平圆形}：恒定速率转弯，半径 5000 m 和 3000 m。
  \item \textbf{垂直弧线 60°、90°}：恒定航向下俯仰渐变拉起。
  \item \textbf{半筋斗保持（60°--150°）}：各相位角的部分筋斗。
\end{itemize}

\subsection{评估指标}

我们报告完成率、存活率、CTE（均值、P90、最大值）、速度切线误差、机头切线误差、机头速度误差、翼面误差、四元数测地误差、攻角/侧滑包线、空速包线、过载、高度包线、能量保持，以及对环类机动的相位分段几何误差。

%===============================================================================

\section{实验结果}

我们将结果按七个研究问题组织。

\subsection{Q1：无 PID 的四元数 RL 技能能否跟踪姿态--速度目标？}

在 2000 步 episode 的高保真评估中，基准四元数技能达到 85\% 存活率，而低保真训练策略在同一环境中的存活率为 0\%（表~\ref{tab:survival}）。
四元数策略的均值航向跟踪误差为 $49.3^\circ$（全 episode RMSE，包含瞬态），过载违规率每步 0.23\%。

\begin{table}[t]
\caption{高保真评估中的存活率与安全性，2000 步 episode。}
\label{tab:survival}
\centering
\begin{tabular}{lccc}
  \hline
  \textbf{策略} & \textbf{存活率} & \textbf{平均 Ep. 时长 (s)} & \textbf{G 违规率} \\
  \hline
  低保真策略，高保真环境  & 0.00  & 2.1   & 1.40\% \\
  高保真策略，高保真环境   & 0.85  & 340.2 & 0.23\% \\
  \hline
\end{tabular}
\end{table}

该技能在 78\% 的目标切换中，于允许的检查间隔内收敛到 $7^\circ$ 四元数误差以内。

\subsection{Q2：为什么四元数目标编码优于欧拉目标？}

对于中等机动（俯仰 $< 45^\circ$），欧拉和四元数编码性能相当。
当俯仰接近 $\pm90^\circ$ 时，由于运动学奇异性（航向在此处无定义），欧拉表示性能退化。
四元数误差 $q_{\text{err}} = q_{\text{tgt}} \otimes q_{\text{curr}}^*$ 在所有姿态下保持良好定义，在 SO(3) 的切空间中提供平滑的 3 自由度误差信号。
对于涉及倒飞的机动，欧拉表示在 $\pm180^\circ$ 处还额外受到滚转不连续性的影响，而四元数表示在整个 SO(3) 流形上是连续的。

\subsection{Q3：目标流组合是否必要？}

我们比较了三种 S 形机动执行方案：（1）静态空间航点，（2）固定 H/P/R/Vt 目标，（3）带参数化前瞻的目标流。
静态航点产生不协调飞行和较大侧滑，因为策略学习到的是被动的"指向航点"策略。
固定 H/P/R/Vt 目标必然失败，因为它们不提供时变姿态信息。
目标流方法对于同时传达\emph{当前期望姿态}和\emph{姿态变化速率}是必要的。

\subsection{Q4：RH-TSO 相比固定目标流参数是否有改进？}

我们将 RH-TSO 与通过网格搜索找到的最佳固定目标流参数化（记为"固定 $\bm{\lambda}$"）进行比较。
表~\ref{tab:rhtso} 报告了四个机动类别的 CTE-P90。

\begin{table}[t]
\caption{RH-TSO 与固定目标流参数对比。$\lambda$ = (前瞻距离 m, $v_t$ m/s)。}
\label{tab:rhtso}
\centering
\begin{tabular}{lcccccc}
  \hline
  \textbf{机动} & \textbf{最佳 $\bm{\lambda}$} & \textbf{固定 CTE-P90} & \textbf{RH-TSO CTE-P90} & \textbf{提升} \\
  \hline
  S 形曲线   & (600, 220)  & 1503 & 918  & 39\% \\
  8 字形     & (600, 220)  & 1039 & 1017 & 2\%  \\
  螺旋线     & (600, 220)  & 691  & 524  & 24\% \\
  90° 垂直弧线 & (1500, 280) & 96   & 51   & 47\% \\
  \hline
\end{tabular}
\end{table}

两个发现尤为突出。
第一，RH-TSO 始终匹配或超过最佳固定参数化，在曲线和垂直机动上提升 24--47\%。
第二，也是更重要的，最佳固定 $\bm{\lambda}$ 因机动类别不同而不同：水平曲线任务（S 形曲线、8 字形、螺旋线）偏好较短的前瞻距离（600 m）和较低的目标速度（220 m/s），而垂直拉起任务偏好较长的前瞻距离（1500 m）和较高的目标速度（280 m/s）。
\emph{不存在适用于所有机动类别的单一固定目标流参数化方案。}
RH-TSO 通过根据当前机动几何和飞机状态在线选择参数来解决这一问题。

\subsection{Q5：能量感知微调是否改善了垂直机动能力？}

基准四元数技能对最高约 $\sim\pm60^\circ$ 的俯仰能可靠跟踪，但超过此点性能退化（表~\ref{tab:vertical}）。
能量感知微调改善了垂直弧线完成率：60° 弧线达到 91\% 完成率（B 级环面质量），90° 弧线达到 82\%（B 级），半筋斗保持最高 150° 达到 68\%（B 级）。
然而，180° 半筋斗仍然是已诊断的失败边界——倒飞顶部阶段（$>150^\circ$）的 VT 误差和翼面误差表明，当前技能无法以可接受的几何质量可靠完成完整筋斗。
我们将 60°--150° 的垂直弧线定性为在能力包线内（环面质量评级 B 级），将 180° 定性为在包线外。

\begin{table}[t]
\caption{基准 vs.\ 能量感知微调技能的垂直机动完成率和环面质量评级。B 级：完成且几何误差中等；F 级：失败或几何不可接受。}
\label{tab:vertical}
\centering
\begin{tabular}{lcccc}
  \hline
  \textbf{机动} & \textbf{基准完成率} & \textbf{微调完成率} & \textbf{环面质量评级} \\
  \hline
  垂直弧线 60°  & 72\% & 91\% & B \\
  垂直弧线 90°  & 41\% & 82\% & B \\
  半筋斗 120°   & 18\% & 68\% & B \\
  半筋斗 150°   & --   & 58\% & B \\
  半筋斗 180°   & --   & $<40\%$ & F（失败边界） \\
  \hline
\end{tabular}
\end{table}

\subsection{Q6：为什么 CTE 单独评估对环形机动会失效？}

CTE 衡量飞机到参考轨迹的垂直距离，忽略了相对机动平面的方向。
一架以偏离预期环面 $15^\circ$ 的翼面旋转飞半筋斗的飞机，可在整个机动中保持 CTE 低于 50 m——然而速度切线误差和翼面误差分别超过 $40^\circ$ 和 $55^\circ$，表明机动本质上是错误的。
相反，一个协调良好的环，即使有中等位置偏差（CTE $\approx$ 120 m），其速度切线误差可低于 $15^\circ$，翼面误差低于 $20^\circ$——这是一个好得多的机动。
我们的几何感知指标集能正确区分这些情况；仅用 CTE 会将它们完全排反。

\subsection{Q7：当前学习到的飞行技能的能力边界在哪里？}

通过系统测试，我们确定了以下能力边界：
\begin{itemize}
  \item \textbf{俯仰}：对于半径 $\geq 3000$ m 的半筋斗，持续拉起至 150° 是可实现的。超过 150°（倒飞到正飞的过渡段），环面几何质量急剧退化。
  \item \textbf{速度}：持续拉起过程中空速低于 150 m/s 导致无法恢复的失速。
  \item \textbf{滚转}：$\pm180^\circ$ 的协调滚转是可执行的，但倒飞阶段超过 $\sim60^\circ$/s 的滚转速率偶尔导致失控。
  \item \textbf{能量}：超过 $\sim\pm8000$ J/kg 的比能变化超出 F-16 推重比包线；策略正确地将油门推至饱和。
  \item \textbf{过渡段}：半筋斗退出过渡段（$>150^\circ$ 相位）是最具挑战性的区域，机头切线误差从 120° 相位时的 $<10^\circ$ 增加到 180° 相位时的 $>25^\circ$。
\end{itemize}

这些边界与仿真 F-16 的物理能力一致，表明学习到的技能内化的是飞机的动态约束，而非记忆特定的机动模板。

%===============================================================================

\section{讨论}

\subsection{为什么优化目标流参数而非作动器？}

在目标流空间中优化而非作动器空间有三个优势：（1）搜索维度低（2--3 个参数 vs.\ 5 个作动器 $\times$ 时域步数），使格点搜索可追踪；（2）每个候选方案都通过冻结策略评估，因此代价曲面尊重技能的实际执行能力；（3）接口可解释——如果 RH-TSO 选择了较长的前瞻距离，我们可以将其解读为"该技能需要为这类机动提前做姿态准备"。

\subsection{为什么选择组合式而非端到端？}

组合式方法将\emph{做什么}（目标流优化）与\emph{怎么做}（学习到的飞行技能）分开。
这反映了四旋翼学习中向将学习组件与结构化中间表示相结合的系统发展的趋势~\citep{kaufmann2023deep,kaufmann2019combining}。
目标流接口是可检查的：我们可以调试失败源自优化器还是飞行技能。

\subsection{局限性}

\begin{enumerate}
  \item \textbf{仿真到现实的鸿沟}：所有结果均在仿真中获得。真实部署需要解决传感器噪声、作动器动态特性、风和建模误差。
  \item \textbf{搜索时域}：当前格点搜索在固定时域上评估有限网格。扩展到带自适应时域的连续优化（CEM/MPPI）是自然的下一步，但需要仔细设计代价函数以避免目标流空间中的局部极小。
  \item \textbf{全包线覆盖}：极端姿态（持续 $>150^\circ$ 俯仰、过失速、尾旋）仍在能力边界之外。180° 半筋斗是已诊断的失败，而非已解决的问题。
  \item \textbf{目标流表达能力}：当前 2--3 个参数的目标流可能不足以应对需要分段速度剖面或非对称俯仰塑形的机动。需要更高维的参数化配合基于采样的优化。
  \item \textbf{多机扩展}：当前框架仅操作单架飞机。
\end{enumerate}

%===============================================================================

\section{结论}

我们提出了一个用于固定翼机动跟踪学习的组合式框架，其核心是一个冻结的、无 PID 的四元数 RL 飞行技能。
我们引入了滚动时域目标流优化（RH-TSO），通过闭环策略 rollout 评估候选方案来在线选择目标流参数，替代了无法跨机动类别泛化的固定启发式参数化方案。
RH-TSO 在曲线和垂直机动上相比固定参数将跟踪精度提高 24--47\%，并揭示了水平曲线任务偏好较短前瞻和较低目标速度、而垂直拉起任务偏好相反——证实了单一固定参数化方案是不够的。
能量感知微调进一步将技能扩展到最高 150° 的垂直弧线（B 级环面质量），180° 半筋斗被定性为已诊断的失败边界。
我们的几何感知评估协议证明，CTE 单独不足以评估环形机动，速度切线、机头切线、翼面和四元数测地误差提供了更忠实的机动质量评估。

未来工作包括将 RH-TSO 扩展到更高维目标流空间并配合基于采样的优化（CEM/MPPI），通过在可执行技能空间中的轨迹优化来自动化目标流生成，以及探索能够弥合固定翼 RL 仿真到现实鸿沟的领域随机化策略。

%===============================================================================

\clearpage
\acknowledgments{作者感谢 Planax 开发团队提供的高保真仿真平台和基线实现。本研究受 [待补充资助信息] 支持。}

%===============================================================================

\bibliography{example}

\end{document}
